\documentclass{article}

% if you need to pass options to natbib, use, e.g.:
%     \PassOptionsToPackage{numbers, compress}{natbib}
% before loading neurips_2026

% The authors should use one of these tracks.
% Before accepting by the NeurIPS conference, select one of the options below.
% 0. "default" for submission
\usepackage{neurips_2026}
% the "default" option is equal to the "main" option, which is used for the Main Track with double-blind reviewing.
% 1. "main" option is used for the Main Track
%  \usepackage[main]{neurips_2026}
% 2. "position" option is used for the Position Paper Track
%  \usepackage[position]{neurips_2026}
% 3. "eandd" option is used for the Evaluations & Datasets Track
 % \usepackage[eandd]{neurips_2026}
% 4. "creativeai" option is used for the Creative AI Track
%  \usepackage[creativeai]{neurips_2026}
% 5. "sglblindworkshop" option is used for the Workshop with single-blind reviewing
 % \usepackage[sglblindworkshop]{neurips_2026}
% 6. "dblblindworkshop" option is used for the Workshop with double-blind reviewing
%  \usepackage[dblblindworkshop]{neurips_2026}
% 7. "education" option is used for the Education Track
%  \usepackage[education]{neurips_2026}


% After being accepted, the authors should add "final" behind the track to compile a camera-ready version.
% 1. Main Track
 % \usepackage[main, final]{neurips_2026}
% 2. Position Paper Track
%  \usepackage[position, final]{neurips_2026}
% 3. Evaluations & Datasets Track
 % \usepackage[eandd, final]{neurips_2026}
% 4. Creative AI Track
%  \usepackage[creativeai, final]{neurips_2026}
% 5. Workshop with single-blind reviewing
%  \usepackage[sglblindworkshop, final]{neurips_2026}
% 6. Workshop with double-blind reviewing
%  \usepackage[dblblindworkshop, final]{neurips_2026}
% Workshop submissions require both a paper title and a workshop title.
% For workshops (5., 6.), the authors should add the name of the workshop, "\workshoptitle" command is used to set the workshop title.
% \workshoptitle{WORKSHOP TITLE}
% 7. Education Track
% \usepackage[education, final]{neurips_2026}

% "preprint" option is used for arXiv or other preprint submissions
 % \usepackage[preprint]{neurips_2026}

% to avoid loading the natbib package, add option nonatbib:
%    \usepackage[nonatbib]{neurips_2026}

\usepackage[utf8]{inputenc} % allow utf-8 input
\usepackage[T1]{fontenc}    % use 8-bit T1 fonts
\usepackage{hyperref}       % hyperlinks
\usepackage{url}            % simple URL typesetting
\usepackage{booktabs}       % professional-quality tables
\usepackage{amsfonts}       % blackboard math symbols
\usepackage{nicefrac}       % compact symbols for 1/2, etc.
\usepackage{microtype}      % microtypography
\usepackage{xcolor}         % colors
\usepackage{amsmath,amssymb,amsthm}
\usepackage{graphicx}
\graphicspath{{../figures/}}
\newtheorem{proposition}{Proposition}
\newtheorem{definition}{Definition}

% Note. For the workshop paper template, both \title{} and \workshoptitle{} are required, with the former indicating the paper title shown in the title and the latter indicating the workshop title displayed in the footnote. 
\title{Score-Space Exposure to Targeted Improvements on Two LLM Leaderboards}

\author{Anonymous Authors}

\begin{document}
\maketitle

\begin{abstract}
How much must a model's task scores improve to become first on a multi-task leaderboard, and how does that depend on aggregation? We study score-space exposure on two public LLM leaderboards, with fixed rivals and independent per-task gain caps. Investment means total score points, tasks touched or largest task gain; it is not measured training effort. We analyse the Open LLM Leaderboard (4495 targets, six tasks; main ordinal results at ranks 2--50) and HELM MMLU (78 targets, 57 subjects). At a $10^{-3}$-point aggregate margin, which is not noise-robust, the mean is feasible for 12.1\% of OLLB targets (95\% pair-bootstrap range 11.2--16.6\%) under marginal caps estimated from 540 fine-tune/base pairs, and 19.2\% of HELM targets under a transferred, uncalibrated allowance; with an oracle 100-point ceiling all targets are feasible. Within the 40-target OLLB stratum at ranks 11--50, paired median costs for median, minimum, Borda and a Condorcet-winner requirement are 0.18--0.50 of the mean, based on 24--40 jointly feasible pairs per comparison. Cardinal closed forms cover the mean, minimum and odd-task median; even-task medians are solved numerically. Ordinal allocations use a finite score grid and bound the continuous strict-win infimum from above. Feasibility and cost need not give the same ordering of rules, and neither establishes an ordering of real training budgets.
\end{abstract}

\section{Introduction}
\label{sec:intro}
Multi-task benchmarks rank models by aggregating per-task scores into one number, and the resulting rank influences which models are adopted, published and funded. A developer therefore has an incentive to improve a chosen target model specifically on the benchmark's sub-datasets, for instance by training on data related to them. How much such targeted improvement is needed to reach first place, and how that amount depends on the aggregation rule, is the question we study on public leaderboard data.

Our approach uses an explicit score-allocation model. We take the score table of a public leaderboard as fixed, let the target raise its own scores task by task up to a cap, and compute the \emph{minimum} raise that makes it strictly first under a given aggregation rule. The cap is the central modelling assumption: it stands for how far training on related data can move one score. We therefore report results under a data-grounded cap and an oracle cap, and report infeasibility (no allowed raise makes the target first) as an outcome rather than a missing value.

\paragraph{Contributions.}
(i) A formulation of the minimum targeted investment with three cost measures ($L_1$, $K$, $L_\infty$) and an explicit family of caps (Section~\ref{sec:setup}). (ii) Score-rule solvers and finite-grid integer programs for rank rules, checked on small instances, with closed forms and proofs for the mean, minimum and odd-task median, and a numerical solver for even-task medians (Section~\ref{sec:solvers}, Appendix~\ref{app:proofs}). (iii) A study of score rules over all non-leading models on two public leaderboards and ordinal rules over explicitly bounded target sets, with bootstrap and evaluation-noise intervals (Sections~\ref{sec:exp}--\ref{sec:results}).

\paragraph{Findings.}
Under the OLLB marginal gain cap and the uncalibrated HELM transfer, first place is out of reach for most models under the three score rules: the feasible fraction is 12\% (mean) on OLLB and 19\% on HELM, and the rule changes how far down the board it is reachable. At OLLB ranks 11--50, paired median cost ratios for the other four main rules are 18\%--50\% of the mean on jointly feasible targets; this is a stratum-specific comparison, not a universal ordering. Two of three working hypotheses are only partly borne out (Section~\ref{sec:results}). All conclusions are conditional on the cap and on the other models being fixed; Section~\ref{sec:disc} lists where they fail.

\section{Related work}
\label{sec:related}
\paragraph{Benchmarks and leaderboards.}
MMLU \citep{hendrycks2021mmlu} aggregates 57 subjects, and HELM \citep{liang2023helm} releases per-scenario scores and raw model outputs for many models; BIG-Bench Hard \citep{suzgun2023bbh} selects hard tasks from BIG-Bench. These provide the multi-task score matrices that we analyse.

\paragraph{Fragility of rankings.}
\citet{dehghani2021lottery} show that the relative order of methods can change with the choice of benchmark tasks (the ``benchmark lottery''). \citet{singh2025illusion} document how private testing and unequal data access distort the Chatbot Arena leaderboard. Both concern what is measured and disclosed; we instead hold the leaderboard fixed and ask what a single developer must change for one model to win.

\paragraph{Aggregation and social choice.}
\citet{lanctot2025vase} treat each task as a voter and use voting theory to build cross-task evaluation; the Handbook of Computational Social Choice \citep{brandt2016handbook} covers the rules (Borda, Condorcet-consistent rules, plurality) we reuse. Our rules are the same families, but we measure the score change needed to make a designated target win.

\paragraph{Training on the test task and contamination.}
\citet{zhou2023cheater} show that benchmark leakage into training can raise scores sharply, and \citet{yang2023rephrased} show that rephrased test samples evade common decontamination. \citet{chen2026incentives} analyse, in a game-theoretic model, the incentives that benchmarks create for strategic post-training. Our analysis is empirical and combinatorial: given a fixed score table, we compute minimum score-rule costs and finite-grid ordinal costs under explicit gain caps, without estimating a cost per training example.

\paragraph{Positioning.}
Closest in question is \citet{chen2026incentives}, who model developers choosing how much to spend on benchmark-specific improvement of a single benchmark score, prove non-existence of equilibria for current protocols and a unique equilibrium for tune-before-test. That analysis concerns one score per model; the choice of aggregation rule over tasks does not enter it. We take the opposite stance: no game and no equilibrium, a fixed table, and the score changes required under each aggregation rule. \citet{colombo2022best} propose a social-choice-based aggregation for NLP benchmarks and show it changes conclusions relative to the mean; we do not propose a rule but measure how costly it is to move a model to first place under each. The computational-social-choice literature studies the analogous \emph{bribery} problem, the cost of changing votes so that a chosen candidate wins under a voting rule \citep{faliszewski2009bribery,brandt2016handbook}. Margin of victory instead asks for the fewest voters whose changed ballots alter an election winner \citep{xia2012margin}. We use a designated target and a score-point cost, rather than unrestricted ballot replacements. The manipulator changes only its own task scores and pays in score points under per-task caps. Our contribution is the cap-conditional comparison on public LLM score tables, with closed forms for score rules and finite-grid programs for rank rules.

\section{Problem setup}
\label{sec:setup}
\subsection{Leaderboard, target and investment}
A leaderboard is a score matrix $X\in[0,100]^{M\times n}$ of $M$ models on $n$ tasks. Fix a target model $i$ with score vector $x=X_{i\cdot}$; all other rows stay fixed. We model a targeted improvement by allowing the developer to raise the target's score on task $t$ to any $y_t\in[x_t,c_t]$, where $c_t$ is a \emph{cap} (Section~\ref{sec:caps}). An aggregation rule $f$ maps the score vectors to a ranking; the target is \emph{first} if it is strictly above every other model under $f$ (by a margin $\varepsilon=10^{-3}$ points for score rules). We report the investment in three ways:
\[
L_1=\sum_t (y_t-x_t),\qquad K=\#\{t: y_t>x_t\},\qquad L_\infty=\max_t (y_t-x_t),
\]
the total points, the number of tasks touched, and the largest single-task jump. The minimum of each is taken separately; for monotone score rules $L_\infty$ is the smallest $d$ for which the cap-saturated uniform lift $y_t=\min(x_t+d,c_t)$ wins. For comparison across rules and datasets we also report $K/n$ and $L_1/n$.

\subsection{Caps: how much can one task be improved?}
\label{sec:caps}
In these data, permitting every task to reach 100 makes every target feasible under the main rules. The cap $c_t$ is therefore the central modelling assumption, and we never collapse it to a single choice. Two caps appear in the main text: \texttt{ceil}, $c_t=100$, an oracle upper bound on what can be done and so a lower bound on any required investment; and \texttt{emp99}, $c_t=\min(100,x_t+g_t^{99})$ where $g_t^{99}$ is the 99th percentile of the observed gains (fine-tune minus its declared base model) on task $t$ in the Open LLM Leaderboard. \texttt{emp99} is the only data-grounded cap. It inherits a selection bias, since only published fine-tunes appear, and it measures gains attained by the community, not the gains a determined developer could attain. Further caps (best existing score, and at most $+5/+10/+20$ points per task) are reported in the appendix. When no $y\le c$ makes the target first, the case is \emph{infeasible}; infeasibility is reported as a rate and never coded as a large investment, and all investment statistics are conditional on feasibility.

\subsection{Aggregation rules}
\label{sec:rules}
Score rules apply a function $f$ to the score vector of each model: the \emph{mean} (the rule of both leaderboards we use), the \emph{median}, and the \emph{minimum} (worst-task score). Rank rules use only the order of models on each task: \emph{Borda} (sum over tasks of the number of models beaten, ties counting one half; equivalent to mean rank) and the \emph{Condorcet-winner} requirement (the target beats every rival on a strict majority of tasks, ties counting one half). Being a Condorcet winner is sufficient for being first under any Condorcet-consistent rule such as Copeland or Schulze, so its investment upper-bounds theirs. Dividing Borda by $n(M-1)$ gives the mean per-task pairwise win rate, a rank aggregation used in benchmark reporting; our downloaded HELM MMLU accuracy aggregate is instead the subject mean (Section~\ref{sec:exp}). Geometric, harmonic and power means, plurality and Copeland are supplementary comparisons. Trimmed-mean code is supplied but no real-data trimmed-mean result is claimed.

\subsection{Minimum-investment solvers}
\label{sec:solvers}
For a score rule let $\tau=\max_{j\neq i}f(X_{j\cdot})$. A target already strictly first has zero cost; otherwise we require $f(y)\ge b=\tau+\varepsilon$. The propositions below give exact finite-margin formulas for this latter case. The even-task median uses the average of the two middle scores and is solved by enumeration. Power means use water filling. For ordinal rules we restrict permitted scores to the original score, rival tie levels $s$, and passing levels $s+\varepsilon$, within the cap. Integer programs optimise on this finite grid; we do not equate their values to an attained minimum of a continuous strict-order problem. Raising the target also reduces rivals' Borda or plurality totals; those changes enter the constraints. HiGHS solves the integer programs with relative optimality tolerance $10^{-4}$ and a 120-second limit; rival constraints are added lazily for Borda and Condorcet. A solved ordinal cost is an optimum on the grid within this tolerance; its feasible allocation upper-bounds the continuous strict-win infimum, rather than establishing equality with it. Checks against small-instance exhaustive search are supplied in \texttt{code/tests/}.

\begin{proposition}[Mean]\label{prop:mean}
Let $f$ be the mean, $h_t=c_t-x_t$ and $D=n(b-f(x))>0$. Feasibility is equivalent to $\sum_t h_t\ge D$. Then $L_1=D$, and $K$ is the least $k$ for which the $k$ largest headrooms sum to at least $D$. The minimum $L_\infty$ is the smallest $d$ with $\sum_t\min(d,h_t)\ge D$; it equals $b-f(x)$ if no cap binds.
\end{proposition}
\begin{proposition}[Minimum]\label{prop:min}
Let $f=\min$. Feasibility is equivalent to $c_t\ge b$ for every $t$. Then $L_1=\sum_t(b-x_t)^+$ and $K=\#\{t:x_t<b\}$.
\end{proposition}
\begin{proposition}[Median, $n$ odd]\label{prop:median}
Let $m=(n+1)/2$, $a=\#\{t:x_t\ge b\}$ and $k=(m-a)^+$. For tasks with $x_t<b$ set $d_t=b-x_t$ if $c_t\ge b$, and $d_t=\infty$ otherwise. Feasibility holds iff at least $k$ finite costs exist. Then $L_1$ is the sum of the $k$ smallest such costs and $K=k$.
\end{proposition}
We verify Propositions~\ref{prop:mean}--\ref{prop:median} numerically against the solvers (Section~\ref{sec:exp}); no theorem is claimed for the rank rules.

\section{Experimental setup}
\label{sec:exp}
\paragraph{Data.}
We use two public leaderboards, each at a fixed, recorded version (checksums in the supplementary reproduction notes). \emph{HELM MMLU}, release v1.13.0 (summary dated 2025-01-10): the exact-match score of 79 models on each of the 57 MMLU subjects ($79\times57$, no missing cell). \emph{Open LLM Leaderboard} (OLLB) \texttt{contents}, dataset commit \texttt{9c09a7c} (2025-03-20): 4576 rows, of which 4497 have distinct model names after we keep the best-scoring variant of each duplicated name, and 4496 after removing one flagged model; each model has six normalised task scores (IFEval, BBH, MATH Lvl 5, GPQA, MUSR, MMLU-PRO), so $n=6$. Merged models are kept in the main analysis and dropped in a robustness slice. The mean we use equals the sites' own aggregate (OLLB \texttt{Average}, HELM ``MMLU All Subjects'') to within $6\times10^{-14}$ points on every model before numerical rounding. OLLB scores are rounded to $10^{-9}$ points at load to remove floating-point artefacts in exact ties; this changes no core feasibility verdict. Neither source declares a licence for the score data; we redistribute only the derived score tables.

\paragraph{Targets, caps and what is reported.}
Every model other than the current leader is a target in turn, with all others fixed (4495 targets on OLLB, 78 on HELM). On OLLB the primary cap is \texttt{emp99} (per-task 99th percentile of the gain of a fine-tune over its declared base model, estimated from 540 such pairs in the same leaderboard). HELM has no base/fine-tune pairs, so its primary cap, \texttt{emp99T}, applies the numerical OLLB MMLU-PRO 99th-percentile gain as a raw-point allowance to all 57 subjects. The related benchmark motivates this anchor, but it is neither a calibrated unit conversion nor a measured HELM response. The oracle cap \texttt{ceil} is reported beside it. All investment statistics are medians over feasible targets, shown together with the feasible fraction. OLLB ordinal rules use finite-grid integer programs, which at $M=4496$ models we run for targets at mean ranks 2--50; we apply a 120\,s limit per instance and report time-outs as censored, never as infeasible.

\section{Results}
\label{sec:results}
All statements below are about targets that can be made first within the cap; the feasible fraction is shown beside every cost, and a rule is called cheaper than another only on targets that are feasible under both.

\paragraph{The closed forms are exact (Q1).}
Propositions~\ref{prop:mean}--\ref{prop:median} agree with the exact solvers on every target of both datasets and every cap in the theory-check grid: the largest deviation in $L_1$ is below $2\times10^{-12}$ points, the deviation in $K$ is zero, and no feasibility verdict differs (Appendix Table~\ref{tab:T2}). The OLLB median is solver-only because $n=6$ is even. Targets already strictly first under a different rule have zero cost; this convention is also applied in the theory check.

\begin{figure}[t]
\centering
\includegraphics[width=\linewidth]{F_ollb_primary.pdf}\\[2pt]
\includegraphics[width=\linewidth]{F_helm_primary.pdf}
\caption{Minimum targeted investment to become first versus the target's current rank under the mean, primary caps. Left: fraction of targets that can reach first at all. Middle and right: tasks touched ($K/n$) and total points per task ($L_1/n$), \emph{among feasible targets}. Top: OLLB, $n=6$, cap \texttt{emp99} (4495 targets; Borda and Condorcet only at ranks 2--50). Bottom: HELM MMLU, $n=57$, cap \texttt{emp99T} (78 targets; the cap is transferred from OLLB). Cost points are bin medians with 95\% target-bootstrap intervals; feasibility points have Wilson intervals (bins double from rank 2; bins with fewer than 3 targets omitted). Generated by \texttt{figures/src/fig\_rank\_rules.py}.}
\label{fig:main}
\end{figure}

\paragraph{Feasibility depends on rank and on the rule (Q3).}
Under the stated caps, permitted score changes cannot make most models first (Figure~\ref{fig:main}, left). On OLLB the feasible fraction under the mean and the median is close to one for targets up to rank a few hundred and near zero beyond rank about one thousand; on HELM every target at ranks 2--10 is feasible under the mean, median, Borda and Condorcet rules, and the fraction falls steeply afterwards (Table~\ref{tab:T3}: for ranks 11--50 only 6 of 40 targets are feasible under the mean and 17 of 40 under the median). The minimum rule is the most restrictive on both datasets, since every task below the best rival's weakest score must be raised and any one cap can block it (Proposition~\ref{prop:min}). The cap, not the rule, decides feasibility far from the top: with the oracle cap \texttt{ceil} every target is feasible (Appendix Figure~\ref{fig:oracle}).

Table~\ref{tab:T4} gives the frontier. On OLLB the fraction of targets that can reach first place under \texttt{emp99} is 12\% under the mean, 10\% under the median and 6\% under the minimum; every target is feasible only up to rank 394 (mean), 90 (median) and 5 (minimum), and the largest feasible rank is 901, 1207 and 2225. On HELM under \texttt{emp99T} the fractions are 19\%, 33\% and 12\%, every target is feasible up to rank 16, 19 and 6, and the largest feasible rank is 16, 29 and 20; the Condorcet requirement reaches rank 26 and Borda 17. On HELM the median is feasible for more targets than the mean. On OLLB it reaches a larger maximum rank but fewer targets overall (451 versus 546); the longest fully feasible prefix is also shorter. These measures of exposure therefore disagree. The largest feasible rank can be an isolated case (OLLB, minimum), which is why the column ``all feasible up to'' is shown.

\paragraph{Uncertainty and sensitivity of the caps.}
Resampling the 540 fine-tune/base gain vectors jointly across tasks (1000 draws, seed 0) and recomputing the caps gives 95\% percentile ranges of 11.2--16.6\%, 8.2--12.2\% and 5.0--6.5\% for OLLB mean/median/minimum feasibility. For the mean, the largest feasible rank ranges from 824 to 1201 and the fully feasible prefix from 346 to 590. These ranges propagate empirical cap estimation, assuming exchangeable pairs; shared bases and publication selection are not accounted for. Setting only the IFEval allowance to zero, while retaining all six tasks and the other allowances, reduces mean feasibility from 546 to 150 targets; zeroing MATH gives 256. The binding allowance depends on the rule: zeroing GPQA reduces minimum-rule feasibility from 273 to 11. Of the 540 observed pairs, 99 gain more than half their respective caps on at least two tasks, none reach all six marginal caps, and 521 lose points on at least one task. This documents the optimism of independent gains; it does not prove a new allocation unattainable.

HELM's 19.2\% is a cap-parameterised result. Transferring the six OLLB quantiles as uniform raw-point allowances of 6.38, 7.06, 12.07, 12.11, 25.70 and 62.99 (GPQA, MMLU-PRO, MuSR, BBH, MATH, IFEval) gives mean-rule shares of 17.9\%, 19.2\%, 46.2\%, 46.2\%, 87.2\% and 100\%, with largest feasible ranks 15, 16, 37, 37, 69 and 79, respectively. The allowances have no HELM-native training calibration, and the OLLB normalised units differ from raw HELM points. Appendix Tables~\ref{tab:capunc}--\ref{tab:helmtransfer} give the cap sensitivity results.

\begin{table}[t]
\caption{Feasibility frontier under the primary cap (\texttt{emp99} on OLLB, \texttt{emp99T} on HELM) and the oracle cap (\texttt{ceil}). \emph{share}: fraction of targets for which first place is reachable (Wilson 95\% intervals); $r^*$: largest rank under the mean among feasible targets (may be an isolated case); \emph{all feas.\ up to}: largest $r$ such that every target of rank $\le r$ is feasible; \emph{first bin $<$ 1/2}: first rank bin (bins double from rank 2) in which fewer than half are feasible. $^\dagger$: targets at ranks 2--50 under the mean only, not comparable with the score rules' fractions. Generated by \texttt{figures/src/tables.py}.}
\label{tab:T4}
\centering\scriptsize
\input{../tables/T4_frontier_ollb.tex}\\[4pt]
\input{../tables/T4_frontier_helm.tex}
\end{table}

\paragraph{Cost, conditional on feasibility, differs strongly by rule (Q2).}
Paired comparisons show lower point costs for the other main rules than for the mean in the displayed strata (Table~\ref{tab:T3}). On OLLB at ranks 11--50 the paired median ratios to the mean are $0.22$ for the median ($n=40$), $0.18$ for the minimum ($n=24$), $0.27$ for Borda ($n=40$) and $0.50$ for the Condorcet-winner requirement ($n=33$). On HELM in the same rank stratum the ratios are $0.09$, $0.04$, $0.55$ and $0.28$, based on 6, 2, 6 and 6 jointly feasible targets, respectively; these small HELM comparisons are indicative only. Separate conditional medians in the table use each rule's feasible population and do not by themselves establish a paired ordering. The reason for lower order-statistic costs is structural: the mean must close the whole summed gap ($n$ times the gap, Proposition~\ref{prop:mean}), whereas order-statistic rules pay only for the cheapest majority of tasks or the weakest tasks. Ordinal costs depend on passing rivals task by task. No uniform ordering is supported across rank bins: at OLLB ranks 2--10, for example, the median-to-mean ratio interval extends above one.

\paragraph{Ordinal rules beyond the top 50.}
A full-board extension draws 15 targets from each mean-rank stratum 51--100, 101--300, 301--1000 and above 1000 (seed 0; 60 targets). Under \texttt{emp99}, Borda is feasible for 15, 15, 7 and 0 targets in these strata; Condorcet for 7, 6, 0 and 0. All 60 are feasible under both rules with \texttt{ceil}; none of the 240 target--rule--cap instances is censored. Under \texttt{emp99}, paired median $L_1$ ratios to the mean are 0.47 [0.42, 0.51] for Borda ($n=35$) and 0.60 [0.55, 0.70] for Condorcet ($n=13$), with 95\% target-bootstrap intervals. Mean feasibility is 36/60, versus 37/60 and 13/60 for the rank rules. Equal sampling by stratum makes these pooled counts and cost ratios descriptions of the sample, not full-board population estimates. The main top-40 cost ratios therefore generalise in direction here, with different magnitudes and selected feasible pairs (Appendix Tables~\ref{tab:ordstatus}--\ref{tab:ordpaired}).

\begin{table}[t]
\caption{Median investment to become first among feasible targets, by rank stratum and rule, primary caps: OLLB (\texttt{emp99}, top) and HELM (\texttt{emp99T}, bottom). ``feas.'' = feasible targets / targets in the stratum; $K$, $L_1$ with 95\% bootstrap intervals over targets; The ratio is the median per-target $L_1$ relative to the mean over targets feasible under both; the final column gives the paired count. Intervals use 2000 target-bootstrap resamples with deterministic data-derived seeds; repeated rows reuse draws. OLLB Borda and Condorcet use targets at ranks 2--50 only. Generated by \texttt{figures/src/tables.py}.}
\label{tab:T3}
\centering\scriptsize
\input{../tables/T3_rules_ollb_emp99.tex}\\[4pt]
\input{../tables/T3_rules_helm_emp99T.tex}
\end{table}

\paragraph{Which of our predictions held.}
Three working hypotheses organize the descriptive comparisons (H1--H3, Appendix~\ref{app:hyp}); we claim no preregistered confirmatory test. H2 is supported as a broad binned trend, not as monotonicity for every target: isolated far-down feasible targets remain, and comparisons across non-nested cap families are not a tightening experiment. H1 (order-statistic and majority rules need \emph{more} tasks but possibly fewer points than the mean) holds only in part: they need fewer points, but not more tasks, except that the Condorcet requirement touches more tasks than the mean on OLLB at $n=6$ ($K=4$ against $2$ at ranks 11--50); on HELM the median needs fewer tasks than the mean ($K=12$ against $46$). H3 is examined next.

\paragraph{Number of tasks (Q4).}
Figure~\ref{fig:scaling} draws random subsets of $n\in\{3,6,12,24,48\}$ HELM subjects and takes the model ranked second or fifth under each rule within each subset as the target. Target identities can therefore differ across rules and subset sizes; this is not a fixed-model causal effect of adding tasks. For the rank-2 target $K/n$ falls with $n$ under every rule, and for the rank-5 target the mean and Borda rules keep touching about $0.4$ of the tasks while the median, minimum and plurality fall. For the rank-5 target the mean rule's $L_1/n$ stays roughly constant in $n$ while that of the median, minimum and plurality falls; for the rank-2 target all rules fall. H3 (that $K/n$ of the median and majority rules does not shrink) is therefore not supported for the median. The bands are wide (60 subset draws), so we make no claim about differences between neighbouring $n$. An OLLB sensitivity replaces six headline tasks with 39 score components on the same 499 models (Appendix~\ref{app:components}); it changes task weighting and does not identify an effect of $n$ alone.

\begin{figure}[t]
\centering
\includegraphics[width=\linewidth]{F_helm_scaling.pdf}
\caption{HELM: investment versus the number of tasks $n$ in a random subset of subjects, for the target ranked second (left) or fifth (right) within the subset under the rule being evaluated; bands span the 2.5--97.5 percentiles over 60 subset draws, among feasible draws, cap \texttt{emp99T}. Generated by \texttt{figures/src/fig\_scaling.py}.}
\label{fig:scaling}
\end{figure}

\paragraph{Robustness (Q5).}
Appendix Figures~\ref{fig:robust}--\ref{fig:robustrank} perturb evaluation scores using independent Gaussian draws with binomial standard errors (300 draws on OLLB score rules, 60 on its ordinal rules, 200 on HELM), and bootstrap tasks with replacement. These are sensitivity distributions, not calibrated confidence intervals for real training costs. The observed $L_1$ lies inside the noise interval in every observed-feasible case examined and usually slightly below the draw median. Near the frontier, even feasibility changes: the HELM mean rank-20 target is infeasible in the observed table but feasible in about 4.5\% of noise draws. Thus a minimum-margin win need not persist under score perturbations. The OLLB noise model transfers item counts from one published run to all models and approximates composite task metrics as binomial; correlations between models and items are unavailable. At a two-point aggregate margin, OLLB mean/median/minimum feasibility falls from 12.1\%/10.0\%/6.1\% to 10.4\%/7.6\%/2.6\%; HELM falls from 19.2\%/33.3\%/11.5\% to 11.5\%/20.5\%/9.0\% (Appendix Table~\ref{tab:margin}). Two points is a sensitivity parameter, not a calibrated confidence threshold. Without shared-item error correlations we cannot infer a re-evaluation success probability from these leads.



\section{Discussion and limitations}
\label{sec:disc}
\paragraph{What the results say.}
The required score investment depends on both the rule and the cap. On the displayed jointly feasible comparisons, order-statistic rules require fewer points than the mean, while their feasibility can be lower: OLLB has 451 feasible targets under the median and 546 under the mean. The median's larger maximum reachable rank is therefore insufficient to call it more vulnerable overall. HELM gives a different feasibility ordering, with 26 median-feasible and 15 mean-feasible targets. Minimum-score aggregation is feasible for few targets under the primary caps, yet costs few points when feasible. There is no single rule ordering that simultaneously describes total gain, feasibility share and the farthest reachable rank. These are score-space requirements under fixed rivals. We do not validate them against held-out fine-tunes, longitudinal leader changes or measured training budgets; calibrating score gains to data, compute or money remains open.

\paragraph{Conditions under which the conclusions fail or are not established.}
(1) \emph{Gain assumptions.} Published fine-tune/base pairs are selection-biased, their upper marginal quantiles are uncertain, and independent improvements need not be jointly attainable. The HELM allowances are arbitrary numerical transfers across different tasks and score units. (2) \emph{Fixed rivals.} Other developers do not respond; strategic interaction needs a game model such as \citet{chen2026incentives}. (3) \emph{Score costs.} Fewer points need not mean less data, compute or money; task-specific or convex training costs could change the ordering. (4) \emph{Evaluation noise.} The primary margin is below score noise; the two-point alternative is a sensitivity parameter. Missing shared-item correlations prevent a re-evaluation probability claim. (5) \emph{Ordinal coverage.} Costs optimise a finite grid, not a continuous strict-win infimum. OLLB main ordinal comparisons cover ranks 2--50; the equal-stratum extension covers 60 further targets, not all 4495. The Condorcet requirement upper-bounds Condorcet-consistent rules; Kemeny and Schulze were not studied. (6) \emph{Scale and snapshots.} HELM has 78 targets; OLLB has six headline tasks. HELM scaling changes target identities and samples one benchmark family; the OLLB component extension changes weighting and includes overlapping GPQA variants. Both snapshots predate later models. (7) \emph{Unresolved costs and unmeasured controls.} Fifty top-500 Copeland point-cost solves are censored; real-data trimmed means and item-level overlap analysis are not supplied.

\paragraph{Broader impact.}
The analysis quantifies how leaderboard rules respond to targeted score improvements under explicit caps. This can inform the choice of aggregation rule and motivate evaluation on held-out tasks, although neither is tested here as a defence against training on benchmark-related data. The allocations could also guide attempts to game rankings. They do not demonstrate that a fine-tune can realise the proposed gains, and we do not train or release an altered model.

\section{Reproducibility}
\label{sec:repro}
Code, pinned public inputs, checksums, environments, seeds, commands and figure/table generators are documented in \texttt{REPRODUCE.md}. The deterministic solvers use no GPU; analysis draws use seed 0 and table bootstraps use data-derived seeds. Core and earlier sensitivity jobs sum to 22,729 wall seconds on the shared two-CPU container; this is neither elapsed time nor CPU-hours. The minimal primary-cap path takes 996 wall seconds and 565 CPU-seconds, excluding retrieval and rendering. The component jobs add 961 wall seconds; cap evidence, wider ordinal jobs and six parallel Copeland shards (8.6--13.3 ks each) are recorded separately, with failures and incomplete accounting disclosed. Data licences are undeclared; we redistribute derived scores only. The forbidden-material record is in \texttt{REPRODUCE.md}.



\bibliographystyle{plainnat}
\bibliography{refs}

\appendix
\section{Proofs of Propositions~\ref{prop:mean}--\ref{prop:median}}
\label{app:proofs}
\paragraph{Proposition~\ref{prop:mean}.} Write $\delta_t=y_t-x_t\in[0,h_t]$. The winning constraint is $\sum_t\delta_t\ge D$. The maximum possible sum is $\sum_t h_t$, proving feasibility. A feasible sum of exactly $D$ is obtained by filling headrooms sequentially and partially filling the last task, so $L_1=D$. Any $k$ touched tasks supply no more than the sum of the $k$ largest headrooms, and filling those tasks attains this bound; this proves the expression for $K$. With maximum raise $d$, no coordinate can supply more than $\min(d,h_t)$, and setting each coordinate to that value attains the bound. Hence the least feasible $d$ gives $L_\infty$. Without binding caps, $nd=D$.
\paragraph{Proposition~\ref{prop:min}.} The condition $\min_t y_t\ge b$ holds iff every coordinate is at least $b$. It is feasible iff every cap is at least $b$. Every deficient coordinate must gain at least $b-x_t$, and setting it exactly to $b$ attains these lower bounds simultaneously. Their sum and count are the two minima.
\paragraph{Proposition~\ref{prop:median}.} For odd $n$, median at least $b$ is equivalent to at least $m$ coordinates at least $b$. The $a$ original coordinates already at least $b$ require no investment. Each other coordinate contributes at most one additional qualifying task, at cost at least $d_t$, attainable exactly when $d_t$ is finite. Thus at least $k$ tasks must be raised, and choosing the $k$ smallest finite costs attains both claimed minima. If fewer than $k$ finite costs exist, feasibility is impossible.

\section{Further results}
\begin{table}[h]
\caption{Closed forms (Propositions~\ref{prop:mean}--\ref{prop:median}) against the exact solvers and an independent generic LP/MILP (60-target sample), over all caps. Generated by \texttt{figures/src/tables.py} from \texttt{results/theory\_check.csv}.}
\label{tab:T2}
\centering\scriptsize
\input{../tables/T2_theory_check.tex}
\end{table}
\begin{figure}[h]
\centering
\includegraphics[width=\linewidth]{F_ollb_oracle.pdf}\\[2pt]
\includegraphics[width=\linewidth]{F_helm_oracle.pdf}
\caption{As Figure~\ref{fig:main} with the oracle cap \texttt{ceil} (scores up to 100): a lower bound on any required investment; every target is feasible, so the differences between rules are not caused by the cap.}
\label{fig:oracle}
\end{figure}
\begin{figure}[h]
\centering
\includegraphics[width=\linewidth]{F_ollb_capgrid.pdf}
\caption{OLLB: effect of the cap on feasibility (left) and cost among feasible targets (right) for three rank strata. Caps are grouped into uniform-gain, empirical-percentile and reference families, which are not ordered against each other. For a fixed target the mean-rule cost is cap-invariant whenever feasible (Proposition~\ref{prop:mean}); conditional medians can change when the feasible target set changes. Generated by \texttt{figures/src/fig\_cap\_grid.py}.}
\label{fig:capgrid}
\end{figure}

\begin{figure}[p]
\centering
\includegraphics[width=\linewidth]{F_helm_capgrid.pdf}
\caption{HELM cap sensitivity, as Appendix Figure~\ref{fig:capgrid}. Empirical percentiles marked T transfer the OLLB MMLU-PRO gain; the pooled percentile transfers gains pooled over six OLLB tasks. Neither is a measured HELM training response. Generated by \texttt{figures/src/fig\_cap\_grid.py}.}
\label{fig:helmcap}
\end{figure}
\begin{figure}[p]
\centering
\includegraphics[width=\linewidth]{F_robust.pdf}
\caption{Score and task sensitivity for selected mean-ranked targets. Panels, in reading order: OLLB evaluation noise, OLLB task bootstrap, HELM evaluation noise, HELM task bootstrap. Open diamonds: observed $L_1$; dots/bars: median and 2.5--97.5 percentiles over feasible draws. Bottom strips show infeasible shares. OLLB score rules use 300 draws; HELM uses 200. Gaussian score perturbations use binomial standard errors, with independent cells; task draws resample columns and their caps together. Generated by \texttt{figures/src/fig\_robust.py}.}
\label{fig:robust}
\end{figure}
\begin{figure}[p]
\centering
\includegraphics[width=\linewidth]{F_robust_rank.pdf}
\caption{Ordinal-rule sensitivity, with the same symbols as Figure~\ref{fig:robust}. OLLB Borda, Condorcet and plurality use 60 noise and task-bootstrap draws at selected targets; HELM contributes 200 evaluation-noise draws only; its ordinal task-bootstrap draws are not plotted. Intervals describe feasible draws and must be read with the infeasible share below. Generated by \texttt{figures/src/fig\_robust.py}.}
\label{fig:robustrank}
\end{figure}

\begin{table}[p]
\centering
\caption{Supplementary rules at the primary caps: OLLB (top; $^\dagger$ plurality uses only targets at ranks 2--50) and HELM (bottom). Conventions and paired-count interpretation follow Table~\ref{tab:T3}; costs and intervals are in score points. Power means shift every score by one point. Generated by \texttt{figures/src/tables.py}.}
\label{tab:supprules}
\input{../tables/T5_supp_rules_ollb.tex}\\[6pt]
\input{../tables/T5_supp_rules_helm.tex}
\end{table}
\begin{table}[p]
\centering
\caption{Larger winning margins (1 and 2 points), primary caps, all targets. Feasible counts have Wilson 95\% intervals; costs are medians over feasible targets (no cost interval in these sensitivity summaries). Base share is feasibility at the primary margin; paired $n$ and ratio use targets feasible with positive cost at both margins. Generated by \texttt{figures/src/tables.py}.}
\label{tab:margin}
\input{../tables/T6a_margin.tex}
\end{table}
\begin{table}[p]
\centering
\caption{OLLB model-set sensitivities. Top: removing merged models (3788 non-leading targets), with a same-target comparison to the main specification. Bottom: restricting targets and rivals to the top 500 models (499 non-leading targets); no full-board comparison is made. Copeland feasibility is established separately from point-cost censoring: $^\ddagger$its cost medians use only the 328 solved-finite targets; 50 $L_1$ solves are censored (46 feasible, four screened infeasible). Feasibility intervals are Wilson 95\%; conditional cost medians are point summaries. Generated by \texttt{figures/src/tables.py}.}
\label{tab:slices}
\input{../tables/T6b_nomerge.tex}\\[6pt]
\input{../tables/T6b_top500.tex}
\end{table}
\begin{table}[p]
\centering
\caption{Sensitivity to score and cost units. Raw OLLB points and task-standard-deviation costs are distinct units, so cross-scale ratios are omitted. Standardisation divides each task's gain by its across-model sample standard deviation. Feasibility intervals are Wilson 95\%; conditional medians are point summaries. Generated by \texttt{figures/src/tables.py}.}
\label{tab:scale}
\input{../tables/T6c_scale.tex}
\end{table}
\begin{table}[p]
\centering
\caption{Weighted means at selected targets, primary caps. Dirichlet concentrations are 1 and 5 (500 weight draws per target); OLLB also includes one-hot weights. $n$ counts target--weight draws, not unique models. The comparator uses uniform weights on the same target. Feasibility intervals are descriptive Wilson intervals over draws; costs are conditional point summaries. Generated by \texttt{figures/src/tables.py}.}
\label{tab:weights}
\input{../tables/T6d_weights.tex}
\end{table}

\clearpage
\section{Reproduction details and interpretation of uncertainty}
\label{app:repro}
The delivery includes raw run CSVs, summaries and generators. For the core results use \texttt{code/run\_core.py}; closed-form checks use \texttt{code/theory\_check.py}; robustness uses \texttt{code/run\_robust.py}; \texttt{code/make\_tables.sh} assembles analysis summaries. Rendering from delivered results uses \texttt{bash figures/src/make\_all.sh results}; that script regenerates every displayed figure and table. A separate source/code/evidence archive is supplied as \texttt{paper/reproducibility.zip}, with per-file SHA-256 hashes inside; regenerate it with \texttt{python3 paper/build\_supplement.py}. Compile this document with \texttt{cd paper \&\& latexmk -pdf main.tex}. The full command sequence and data downloads are in \texttt{REPRODUCE.md} and \texttt{code/run\_all.sh}.

\paragraph{Pinned inputs.} OLLB contents revision \path{9c09a7cae43334062a82cb164f2ef255013dafa2}, parquet SHA-256 \texttt{f13a7061e63133bc8745d89f8c7a5614\allowbreak 8584940ede22d3b27e21cad1e078dca5}; HELM v1.13.0, manifest digest \texttt{830bdb10a5518c50689d60c86623cc3d\allowbreak 79ff5317d7e6b96408ea7dc2058de12b}. OLLB evaluation item counts use one results JSON at revision \path{aa81ecc38fdc5708254b833923368970efdf5ef5}, file SHA-256 \texttt{c0386eb6b6b7ca3c6cc596b190dcd898\allowbreak b17fb7055b092229d00ff50727930856}. Download scripts verify or record checksums; raw downloads are regenerated, while derived matrices are included. Source licences for score data are undeclared; the code's licence does not establish permission for the score data. We use the public numeric scores for research and state this unresolved licence status rather than assign a new licence to them.

\paragraph{Environment and sampling.} Analysis used Python 3.11.2, NumPy 2.4.6, SciPy 1.17.1 (HiGHS), pandas 3.0.6, PyArrow 26.0.0 and pytest 9.2.0; rendering uses Matplotlib 3.11.2 and Pillow 13.0.0. Computation ran on a shared Xeon Gold 6230 container limited to two CPUs and 4 GB RAM. Analysis routines use NumPy's default generator with seed 0; rule-comparison table bootstraps use the deterministic data-derived seeds described below. Each MILP has a 120-second limit; a timeout is censored, never labelled infeasible. Statistics over targets exclude the mean leader and retain zero-cost targets already strictly first under another rule. Thus all counts are indexed by mean rank, not by a separately selected target set for each rule. Ordinal main results use 49 OLLB targets at ranks 2--50, supplemented by the separate 60-target rank-stratified extension, with all 4496 models present (4495 fixed rivals per target); the top-500 sensitivity instead changes the rival set.

\paragraph{Intervals answer different questions.} Figure target-bootstrap intervals and uniform-lift ratios use 1000 resamples; tabulated rule-comparison intervals use 2000 resamples. The latter seed NumPy's default generator from the first eight SHA-256 bytes of the input float vector, interpreted as a little-endian integer; identical vectors reuse draws regardless of table or call order. Thus repeated mean-reference rows agree between Tables~\ref{tab:T3} and~\ref{tab:supprules}. These intervals describe heterogeneity within the chosen target stratum. Wilson intervals accompany displayed feasibility proportions. These score tables enumerate a fixed population; those intervals are descriptive generalisations under a target-exchangeability model, not uncertainty in the observed count itself. Task bootstrap samples $n$ columns with replacement; task-subset scaling samples without replacement. Noise draws independently perturb every cell by a Gaussian with plug-in binomial standard error, clipped to $[0,100]$; this omits shared-item error correlation and approximates composite OLLB metrics. Noise/task bands are percentile spreads, not confidence intervals for the median. None measures variability of real fine-tuning.

\paragraph{Empirical gain magnitudes.} The OLLB cap uses 540 matched fine-tune/base pairs in the frozen table. In task order IFEval, BBH, MATH Lvl 5, GPQA, MUSR, MMLU-PRO, the 99th-percentile headrooms are approximately 62.988, 12.109, 25.696, 6.376, 12.069 and 7.064 points, as generated by \texttt{empirical\_gain\_caps(99)} in \texttt{code/mti/data.py}. HELM transfers the final value to each subject. These are observed marginal quantiles, not fitted returns per training example; the particularly large IFEval allowance illustrates the heterogeneity obscured by a uniform-gain cap.

\paragraph{Capped uniform improvement as a baseline.} The feasible uniform lift for score rules is $y_t=\min(x_t+d,c_t)$, where $d=L_\infty$; its actual point cost is $\sum_t\min(d,c_t-x_t)$. For the mean this equals the minimum $L_1$ by Proposition~\ref{prop:mean}, while touching more tasks. Among positive-cost feasible targets at the primary cap, its median cost ratio to optimal targeted allocation is 2.38 [2.36, 2.42] for the OLLB median ($n=450$ targets) and 3.64 [3.58, 3.74] for the OLLB minimum ($n=272$); HELM gives 8.96 [8.13, 11.35] ($n=26$) and 37.54 [23.11, 53.75] ($n=8$), respectively. These are paired ratios with 95\% target-bootstrap intervals from \texttt{uniform\_vs\_optimal\_*.csv}; zero-cost already-first targets are excluded. Concentrating score gains therefore saves points under these order-statistic rules in this score-space model.

\paragraph{Ordinal numerical scope.} Reducing the crossing grid from $10^{-3}$ to $10^{-5}$ points changes no feasibility verdict in the examined HELM and OLLB top-50 blocks (\texttt{eps\_grid\_check.csv}). One OLLB Borda target changes its minimum number of tasks; the largest cost change is 0.1055 points. We retain the finite-grid definition rather than claim equivalence to the continuous strict-win infimum. The cap can lie between a rival score and its passing level, and numerical optimality tolerances also matter. Ordinal $L_\infty$ in the code is a bisection over strict rank comparisons (tolerance $10^{-6}$), whereas ordinal $L_1$ and $K$ use the discrete grid; we use score-rule baseline ratios above to keep the winning convention aligned. The rounding reruns change no core feasibility verdict, and all core top-50 $L_1$ changes are below $3\times10^{-9}$ points. The ordinal task-bootstrap rerun resolves one previously censored instance; superseded attempts are recorded in \texttt{REPRODUCE.md}.

\paragraph{Supplementary rules and sensitivities.} The geometric, harmonic and $p=-2$ power means use a one-point shift to handle zeros: $f_p(x)=(n^{-1}\sum_t(x_t+1)^p)^{1/p}-1$, with the geometric limit at $p=0$. Plurality sums first-place task credits, dividing credit equally among tied leaders; the implemented strict-win credit advantage is $10^{-4}$. Copeland counts pairwise task-majority wins with one-half credit for pairwise ties. The Condorcet requirement is stronger: it requires every pairwise contest to be won. These supplementary definitions use the same score caps as the main rules (Table~\ref{tab:supprules}). The mean is not always the most costly score rule in this expanded set: at HELM ranks 11--50 geometric and harmonic means have paired median point-cost ratios of 1.04 and 1.08 relative to the mean, with 6 and 4 jointly feasible targets respectively.

\paragraph{Changing the margin and score scale.} Tables~\ref{tab:margin}--\ref{tab:weights} report the sensitivity summaries. The primary $10^{-3}$-point margin defines a mathematical win, not a noise-robust win. With a two-point aggregate margin, feasible shares under mean/median/minimum fall to 10.4\%/7.6\%/2.6\% on OLLB and 11.5\%/20.5\%/9.0\% on HELM (\texttt{summary\_variants.csv}). Removing merged OLLB models gives 11.5\%/9.4\%/7.0\% under the same empirical cap; these are a different target and rival population. Under a uniform ten-point cap on \emph{raw} OLLB accuracies, the mean, median and minimum feasible shares are 6.6\%, 0.47\% and 5.0\%. A ten-point raw gain is a different intervention from ten normalised points, and costs across these scales should not be interpreted as comparable training effort. These observed population proportions have no sampling error for the fixed snapshots.

\paragraph{Top-500 Copeland.} The supplementary OLLB slice includes Copeland at \texttt{emp99}, with 500 models as targets and rivals (Table~\ref{tab:slices}). Of the 499 targets other than the mean leader, 328 have a solved finite minimum $L_1$, 121 have an infeasible $L_1$ solve, and 50 point-cost solves are censored at the 120-second limit. Separate $K$ solutions establish feasibility for 46 of the censored targets; the other four have infinite $K$ and lose at their cap-saturated vector, certifying infeasibility. Feasibility is therefore settled at 374/499, while minimum point costs are solved for only 328 feasible targets. Their conditional median $L_1$ is 38.1 points and median $K$ is 4; selection by solver completion can affect these summaries. All 500 feasibility verdicts from the separate cost solves agree with the cap-saturated screen; winning vectors establish feasibility, not optimal point costs. Copeland at the oracle cap was not run on this slice. These results are supplementary; the slice changes the rival set as well as the targets.

\begin{table}[ht]
\centering
\caption{Copeland at \texttt{emp99} on the top-500 OLLB slice, by original mean rank. Certified feasibility uses the separate $K$ solves and cap-saturation check; no target's feasibility is unresolved. ``$L_1$ solved'' counts finite optimal point costs; infeasible cases have no finite cost. The 50 point-cost timeouts split into 46 feasible targets with unknown minimum $L_1$ and four certified infeasible targets. $^\ddagger$Both cost medians use only the solved-finite subset, so they can be selected by solver difficulty. Wilson intervals describe the certified fraction under the target-exchangeability interpretation of Appendix~\ref{app:repro}. No oracle Copeland run is included. Generated by \texttt{figures/src/tables.py}.}
\label{tab:copeland}
\input{../tables/T6b_copeland.tex}
\end{table}

\paragraph{Other controls.} The wider ordinal extension instead retains all 4496 models. Its memory-lean MILPs construct rival rows on demand and are checked against the original dense solvers; a cap-saturated screen certifies infeasibility before optimisation. All 240 instances resolve, with no grid-only infeasibility or censoring. Its target-bootstrap ratios use 1000 draws, seed 0; equal-stratum samples are not weighted to the full model population. A timeout is an unresolved instance, not evidence of infeasibility; the rounded-score Borda reruns resolve the previously censored target at each cap. Weighted means draw 500 vectors each from symmetric Dirichlet distributions with concentrations 1 and 5, at seven selected OLLB targets and six HELM targets, with uniform weights as the comparator (OLLB also includes one-hot weights). These observations count target--weight draws, not independent models. Standardised-cost runs minimise $\sum_t(y_t-x_t)/\sigma_t$, where $\sigma_t$ is the across-model sample standard deviation on task $t$; they change the cost function while retaining the rule and caps. OLLB's even-task median is not included in that cost control. No standardised score cost is equated to actual training resources.

\paragraph{Traceability.} Numeric input hashes for each plot are in its adjacent JSON; table input hashes are in \texttt{tables/MANIFEST.json}. Theory checks trace to \texttt{results/theory\_check.csv}; feasibility and paired comparisons trace to \texttt{results/core\_*.csv} and \texttt{results/summary/}. The single forbidden-material encounter record is in \texttt{REPRODUCE.md}, section ``Forbidden-material record''; none is reported.

\section{Cap sensitivity and wider ordinal coverage}
\label{app:capevidence}
Tables~\ref{tab:capunc}--\ref{tab:allowance} propagate cap-estimation uncertainty from the 540 observed gain vectors. They leave the six-task leaderboard fixed. The pair bootstrap assumes exchangeable pairs; shared base models, publication selection and the choice of the 99th percentile remain outside these intervals. The observed gains need not jointly attain the marginal caps. Table~\ref{tab:helmtransfer} changes the numerical allowance on HELM without claiming an empirical HELM training response.

\begin{table}[ht]
\centering
\caption{OLLB feasibility under the observed \texttt{emp99} caps, with 2.5--97.5 percentile ranges over 1000 joint pair-bootstrap draws (seed 0). The denominator is all 4495 non-leading targets. Ranks are under the original mean; the prefix ends at the largest rank for which every target from rank 2 is feasible. Ranges concern estimation of the caps from these pairs, rather than training uncertainty or sampling of target models. Generated by \texttt{figures/src/tables.py}.}
\label{tab:capunc}
\input{../tables/T8a_cap_uncertainty.tex}
\end{table}

\begin{table}[ht]
\centering
\caption{OLLB allowance ablation. Each row sets the named task's gain allowance to zero, retaining all six tasks and the other five allowances. Feasible counts out of 4495 targets and the omitted task's original cap are observed values [2.5, 97.5 percentiles over the same 1000 pair-bootstrap draws]. The primary row repeats Table~\ref{tab:capunc}. Generated by \texttt{figures/src/tables.py}.}
\label{tab:allowance}
\input{../tables/T8b_allowance_ablation.tex}
\end{table}

\begin{table}[ht]
\centering
\caption{HELM allowance sensitivity on 78 targets. The first six rows transfer the named OLLB task's \texttt{emp99} value as a uniform raw-point allowance $h$ on all 57 subjects; $^\star$ marks the primary \texttt{emp99T} cap. These are uncalibrated transfers between different score units. MuSR and BBH differ before rounding (12.07 and 12.11), but give identical counts. Remaining rows use uniform allowances of 5, 10 and 20 points. Cells show feasible count (percentage); the mean-rule share also has a 95\% target-bootstrap interval. Prefix and maximum rank follow Table~\ref{tab:capunc}. Generated by \texttt{figures/src/tables.py}.}
\label{tab:helmtransfer}
\input{../tables/T9_helm_allowance.tex}
\end{table}

\clearpage
Tables~\ref{tab:ordstatus}--\ref{tab:ordpaired} examine 60 further OLLB targets with all 4496 models retained as rivals or target. Fifteen targets are sampled from each of four mean-rank strata (seed 0, realised ranks 51--4104). Equal allocation to strata of different sizes makes the pooled counts descriptive of this sample, rather than estimates of full-board feasibility. A losing cap-saturated vector certifies infeasibility. Every remaining finite-grid instance is solved within the stated numerical tolerance; no target is censored or found infeasible by the MILP.

\begin{table}[ht]
\centering
\caption{Wider OLLB ordinal coverage: feasible counts out of 15 targets per rank stratum and out of 60 overall. The final column separates cap-screen infeasibility from solver-limit censoring. Neither rule has any censored instance or already-first target. The oracle cap allows scores up to 100. Generated by \texttt{figures/src/tables.py}.}
\label{tab:ordstatus}
\input{../tables/T10a_ordinal_sample.tex}
\end{table}

\begin{table}[ht]
\centering
\caption{Paired ordinal costs on the same 60-target sample. Rule and mean feasibility counts precede the number feasible under both, which is the denominator for the median per-target cost ratios and their 95\% target-bootstrap intervals (1000 draws, seed 0). These selected pairs are 35 targets for Borda and 13 for Condorcet under \texttt{emp99}; ratios do not summarise infeasible targets. Generated by \texttt{figures/src/tables.py}.}
\label{tab:ordpaired}
\input{../tables/T10b_ordinal_paired.tex}
\end{table}

\clearpage
\section{OLLB score-component sensitivity}
\label{app:components}
We use the pinned OLLB per-model results snapshot to replace the six headline tasks with 39 score components: 24 BBH, three GPQA, seven MATH-hard and three MuSR scores, plus MMLU-PRO and IFEval. The 500 models with the highest six-task means are the retrieval population. For 499 of them, a result file matches all six published raw task scores to $10^{-6}$; one model at original rank 432 has no matching file and is excluded without imputation. Among matching files we choose the latest. Scores subtract the task's chance baseline, rescale to $[0,100]$ and clip at zero, as documented in \texttt{results/leaf/README.md}. The comparator uses the six published task scores of these same 499 models. This is a change in score representation and weighting: BBH contributes 24/39 of the component mean, versus 1/6 of the headline mean. These data therefore do not isolate the effect of increasing task count, and the components are not assumed independent. In particular, the GPQA diamond/main/extended entries are overlapping dataset variants, not disjoint item sets; item-level overlap was not audited.

The score rules are run on all 499 models, including each representation's current leader; paired summaries exclude the six-task mean leader, leaving 498 targets. Rank rules use the top 50 under each representation's own mean; their same-model intersection after excluding the six-task mean leader contains 40 targets. The oracle and ten-point caps are applied in each representation. The empirical analogue transfers each headline task's 99th-percentile gain to all of its components, rather than estimating a component-specific training response. Only 23 declared fine-tune/base pairs occur in this model slice, so component-specific empirical percentiles are not estimated. Neither component noise nor component-resampling intervals, nor component Copeland results, are claimed.

Table~\ref{tab:components} reports the same-model comparisons. Under the transferred empirical caps, mean feasibility is 464/498 in the component table and 472/498 in the six-task table; median feasibility is 220/498 versus 396/498, while minimum feasibility is 242/498 versus 201/498 (\texttt{paired\_leaf\_vs\_six.csv}). Thus the change in representation can reverse the direction of a feasibility difference across rules. On jointly feasible, positive-cost targets, the paired median ratio of points per component to points per headline task is approximately 0.84 for the mean ($n=461$ pairs), 0.39 for the median ($n=201$) and 0.40 for the minimum ($n=94$). These ratios normalise total costs by 39 and 6, respectively; their target-bootstrap intervals describe heterogeneity across the chosen slice, not uncertainty in fine-tuning returns. For positive-cost oracle-cap pairs, the median ratio of $K/39$ to $K/6$ is 0.62 for the mean, 0.67 for the median and 0.38 for the minimum (paired counts 497, 496 and 496); under transferred empirical caps the mean ratio is 1.02, while the median ratio is 0.51. Thus lower task fractions are not universal across cap choices. All component runs completed without censoring. The gate, selected file names, derived matrix, individual runs and summary generators are supplied under \texttt{data/processed/ollb\_leaves.csv}, \texttt{results/ollb\_leaves\_gate.json} and \texttt{results/leaf/}. The source manifest and selected-file manifest there identify 1018 retrieved files and their SHA-256 hashes; \texttt{source\_digests.json} records the reduced-retrieval digest and derived matrix digest.

\begin{table}[p]
\centering
\caption{OLLB score-component sensitivity on the same 499-model rival set. The score-rule panel examines all 498 targets except the six-task mean leader; $^\dagger$ ordinal rows examine the 40 models in both top-50 target sets, after the same exclusion. Feasible counts are within each examined set; cost ratios and their 95\% target-bootstrap intervals (1000 resamples, seed 0) use only targets feasible with positive cost in both representations. $^{\S}$The empirical cap is transferred from each parent task to its components. The 39-component mean weights BBH as 24/39, versus 1/6 in the six-task mean; this is not a pure task-count comparison. Nothing is censored. Generated by \texttt{figures/src/tables.py}.}
\label{tab:components}
\input{../tables/T7_leaf_vs_six.tex}
\end{table}

\section{Working hypotheses and their interpretation}
\label{app:hyp}
These working hypotheses are motivated by the closed forms and organize descriptive comparisons; no preregistration is claimed. H1: with loose caps, score-sum rules need the same $L_1$ regardless of $K$, while order-statistic and majority rules need more tasks but possibly fewer points than the mean. H2: required investment grows with the target's initial rank; feasibility falls as caps tighten; rule ordering can change with the cap. H3: for fixed per-task caps the fraction $K/n$ of tasks that majority and median rules need does not shrink with $n$, while the mean rule's $L_1/n$ depends on the density of gaps. The analysis is exploratory. Lower task counts contradict the task-concentration part of H1, as observed; higher paired point costs would contradict its cost comparison for the relevant rule and stratum. A falling median $K/n$ contradicts H3, as observed. H2 is evaluated as a binned rank trend, not a claim of target-wise monotonicity; feasibility monotonicity for nested caps follows from the model.


\clearpage
\section*{NeurIPS Paper Checklist}




\begin{enumerate}

\item {\bf Claims}
    \item[] Question: Do the main claims made in the abstract and introduction accurately reflect the paper's contributions and scope?
    \item[] Answer: \answerYes{}
    \item[] Justification: Abstract and introduction state the two dataset-specific findings, the cap assumption and their scope; each is backed by Table~\ref{tab:T3}, Table~\ref{tab:T4} and Figure~\ref{fig:main}, and Section~\ref{sec:disc} lists where they fail.
    \item[] Guidelines:
    \begin{itemize}
        \item The answer \answerNA{} means that the abstract and introduction do not include the claims made in the paper.
        \item The abstract and/or introduction should clearly state the claims made, including the contributions made in the paper and important assumptions and limitations. A \answerNo{} or \answerNA{} answer to this question will not be perceived well by the reviewers. 
        \item The claims made should match theoretical and experimental results, and reflect how much the results can be expected to generalize to other settings. 
        \item It is fine to include aspirational goals as motivation as long as it is clear that these goals are not attained by the paper. 
    \end{itemize}

\item {\bf Limitations}
    \item[] Question: Does the paper discuss the limitations of the work performed by the authors?
    \item[] Answer: \answerYes{}
    \item[] Justification: Section~\ref{sec:disc} (Discussion and limitations) lists cap selection bias and transfer, task independence, fixed rivals, noise, rule coverage, snapshots and unfinished items.
    \item[] Guidelines:
    \begin{itemize}
        \item The answer \answerNA{} means that the paper has no limitation while the answer \answerNo{} means that the paper has limitations, but those are not discussed in the paper. 
        \item The authors are encouraged to create a separate ``Limitations'' section in their paper.
        \item The paper should point out any strong assumptions and how robust the results are to violations of these assumptions (e.g., independence assumptions, noiseless settings, model well-specification, asymptotic approximations only holding locally). The authors should reflect on how these assumptions might be violated in practice and what the implications would be.
        \item The authors should reflect on the scope of the claims made, e.g., if the approach was only tested on a few datasets or with a few runs. In general, empirical results often depend on implicit assumptions, which should be articulated.
        \item The authors should reflect on the factors that influence the performance of the approach. For example, a facial recognition algorithm may perform poorly when image resolution is low or images are taken in low lighting. Or a speech-to-text system might not be used reliably to provide closed captions for online lectures because it fails to handle technical jargon.
        \item The authors should discuss the computational efficiency of the proposed algorithms and how they scale with dataset size.
        \item If applicable, the authors should discuss possible limitations of their approach to address problems of privacy and fairness.
        \item While the authors might fear that complete honesty about limitations might be used by reviewers as grounds for rejection, a worse outcome might be that reviewers discover limitations that aren't acknowledged in the paper. The authors should use their best judgment and recognize that individual actions in favor of transparency play an important role in developing norms that preserve the integrity of the community. Reviewers will be specifically instructed to not penalize honesty concerning limitations.
    \end{itemize}

\item {\bf Theory assumptions and proofs}
    \item[] Question: For each theoretical result, does the paper provide the full set of assumptions and a complete (and correct) proof?
    \item[] Answer: \answerYes{}
    \item[] Justification: Propositions~\ref{prop:mean}--\ref{prop:median} state their assumptions in Section~\ref{sec:solvers} and are proved in Appendix~\ref{app:proofs}; rank-rule results are computational and no theorem is claimed.
    \item[] Guidelines:
    \begin{itemize}
        \item The answer \answerNA{} means that the paper does not include theoretical results. 
        \item All the theorems, formulas, and proofs in the paper should be numbered and cross-referenced.
        \item All assumptions should be clearly stated or referenced in the statement of any theorems.
        \item The proofs can either appear in the main paper or the supplemental material, but if they appear in the supplemental material, the authors are encouraged to provide a short proof sketch to provide intuition. 
        \item Inversely, any informal proof provided in the core of the paper should be complemented by formal proofs provided in appendix or supplemental material.
        \item Theorems and Lemmas that the proof relies upon should be properly referenced. 
    \end{itemize}

    \item {\bf Experimental result reproducibility}
    \item[] Question: Does the paper fully disclose all the information needed to reproduce the main experimental results of the paper to the extent that it affects the main claims and/or conclusions of the paper (regardless of whether the code and data are provided or not)?
    \item[] Answer: \answerYes{}
    \item[] Justification: Section~\ref{sec:setup}--\ref{sec:exp} define rules, caps, margin and filters; \texttt{REPRODUCE.md} in the supplementary code gives data versions, seeds and commands.
    \item[] Guidelines:
    \begin{itemize}
        \item The answer \answerNA{} means that the paper does not include experiments.
        \item If the paper includes experiments, a \answerNo{} answer to this question will not be perceived well by the reviewers: Making the paper reproducible is important, regardless of whether the code and data are provided or not.
        \item If the contribution is a dataset and\slash or model, the authors should describe the steps taken to make their results reproducible or verifiable. 
        \item Depending on the contribution, reproducibility can be accomplished in various ways. For example, if the contribution is a novel architecture, describing the architecture fully might suffice, or if the contribution is a specific model and empirical evaluation, it may be necessary to either make it possible for others to replicate the model with the same dataset, or provide access to the model. In general. releasing code and data is often one good way to accomplish this, but reproducibility can also be provided via detailed instructions for how to replicate the results, access to a hosted model (e.g., in the case of a large language model), releasing of a model checkpoint, or other means that are appropriate to the research performed.
        \item While NeurIPS does not require releasing code, the conference does require all submissions to provide some reasonable avenue for reproducibility, which may depend on the nature of the contribution. For example
        \begin{enumerate}
            \item If the contribution is primarily a new algorithm, the paper should make it clear how to reproduce that algorithm.
            \item If the contribution is primarily a new model architecture, the paper should describe the architecture clearly and fully.
            \item If the contribution is a new model (e.g., a large language model), then there should either be a way to access this model for reproducing the results or a way to reproduce the model (e.g., with an open-source dataset or instructions for how to construct the dataset).
            \item We recognize that reproducibility may be tricky in some cases, in which case authors are welcome to describe the particular way they provide for reproducibility. In the case of closed-source models, it may be that access to the model is limited in some way (e.g., to registered users), but it should be possible for other researchers to have some path to reproducing or verifying the results.
        \end{enumerate}
    \end{itemize}


\item {\bf Open access to data and code}
    \item[] Question: Does the paper provide open access to the data and code, with sufficient instructions to faithfully reproduce the main experimental results, as described in supplemental material?
    \item[] Answer: \answerYes{}
    \item[] Justification: The data are public (Open LLM Leaderboard and HELM, versions and checksums in \texttt{REPRODUCE.md}); the supplementary code reproduces every table and figure with the commands listed there.
    \item[] Guidelines:
    \begin{itemize}
        \item The answer \answerNA{} means that paper does not include experiments requiring code.
        \item Please see the NeurIPS code and data submission guidelines (\url{https://neurips.cc/public/guides/CodeSubmissionPolicy}) for more details.
        \item While we encourage the release of code and data, we understand that this might not be possible, so \answerNo{} is an acceptable answer. Papers cannot be rejected simply for not including code, unless this is central to the contribution (e.g., for a new open-source benchmark).
        \item The instructions should contain the exact command and environment needed to run to reproduce the results. See the NeurIPS code and data submission guidelines (\url{https://neurips.cc/public/guides/CodeSubmissionPolicy}) for more details.
        \item The authors should provide instructions on data access and preparation, including how to access the raw data, preprocessed data, intermediate data, and generated data, etc.
        \item The authors should provide scripts to reproduce all experimental results for the new proposed method and baselines. If only a subset of experiments are reproducible, they should state which ones are omitted from the script and why.
        \item At submission time, to preserve anonymity, the authors should release anonymized versions (if applicable).
        \item Providing as much information as possible in supplemental material (appended to the paper) is recommended, but including URLs to data and code is permitted.
    \end{itemize}


\item {\bf Experimental setting/details}
    \item[] Question: Does the paper specify all the training and test details (e.g., data splits, hyperparameters, how they were chosen, type of optimizer) necessary to understand the results?
    \item[] Answer: \answerYes{}
    \item[] Justification: Section~\ref{sec:exp} and Appendix state data filters, dedupe, caps, bootstrap and Monte Carlo settings; \texttt{REPRODUCE.md} lists the rest.
    \item[] Guidelines:
    \begin{itemize}
        \item The answer \answerNA{} means that the paper does not include experiments.
        \item The experimental setting should be presented in the core of the paper to a level of detail that is necessary to appreciate the results and make sense of them.
        \item The full details can be provided either with the code, in appendix, or as supplemental material.
    \end{itemize}

\item {\bf Experiment statistical significance}
    \item[] Question: Does the paper report error bars suitably and correctly defined or other appropriate information about the statistical significance of the experiments?
    \item[] Answer: \answerYes{}
    \item[] Justification: Target-bootstrap cost intervals and Wilson feasibility intervals (Figure~\ref{fig:main}, Tables~\ref{tab:T3}, \ref{tab:T4}), cap-estimation pair-bootstrap ranges (Section~\ref{sec:results}), plus noise/task percentile spreads on both datasets (Figures~\ref{fig:robust}--\ref{fig:robustrank}); assumptions and their scope are in Appendix~\ref{app:repro}.
    \item[] Guidelines:
    \begin{itemize}
        \item The answer \answerNA{} means that the paper does not include experiments.
        \item The authors should answer \answerYes{} if the results are accompanied by error bars, confidence intervals, or statistical significance tests, at least for the experiments that support the main claims of the paper.
        \item The factors of variability that the error bars are capturing should be clearly stated (for example, train/test split, initialization, random drawing of some parameter, or overall run with given experimental conditions).
        \item The method for calculating the error bars should be explained (closed form formula, call to a library function, bootstrap, etc.)
        \item The assumptions made should be given (e.g., Normally distributed errors).
        \item It should be clear whether the error bar is the standard deviation or the standard error of the mean.
        \item It is OK to report 1-sigma error bars, but one should state it. The authors should preferably report a 2-sigma error bar than state that they have a 96\% CI, if the hypothesis of Normality of errors is not verified.
        \item For asymmetric distributions, the authors should be careful not to show in tables or figures symmetric error bars that would yield results that are out of range (e.g., negative error rates).
        \item If error bars are reported in tables or plots, the authors should explain in the text how they were calculated and reference the corresponding figures or tables in the text.
    \end{itemize}

\item {\bf Experiments compute resources}
    \item[] Question: For each experiment, does the paper provide sufficient information on the computer resources (type of compute workers, memory, time of execution) needed to reproduce the experiments?
    \item[] Answer: \answerYes{}
    \item[] Justification: Section~\ref{sec:repro} and \texttt{REPRODUCE.md} give the CPU type, memory and measured job wall-clock durations; several timings and CPU-seconds are unavailable, and no compliance with the unmeasured CPU-hour estimate is asserted.
    \item[] Guidelines:
    \begin{itemize}
        \item The answer \answerNA{} means that the paper does not include experiments.
        \item The paper should indicate the type of compute workers CPU or GPU, internal cluster, or cloud provider, including relevant memory and storage.
        \item The paper should provide the amount of compute required for each of the individual experimental runs as well as estimate the total compute. 
        \item The paper should disclose whether the full research project required more compute than the experiments reported in the paper (e.g., preliminary or failed experiments that didn't make it into the paper). 
    \end{itemize}
    
\item {\bf Code of ethics}
    \item[] Question: Does the research conducted in the paper conform, in every respect, with the NeurIPS Code of Ethics \url{https://neurips.cc/public/EthicsGuidelines}?
    \item[] Answer: \answerYes{}
    \item[] Justification: The work uses only public score tables, trains no models and involves no human subjects; we have read the NeurIPS Code of Ethics and see no conflict.
    \item[] Guidelines:
    \begin{itemize}
        \item The answer \answerNA{} means that the authors have not reviewed the NeurIPS Code of Ethics.
        \item If the authors answer \answerNo, they should explain the special circumstances that require a deviation from the Code of Ethics.
        \item The authors should make sure to preserve anonymity (e.g., if there is a special consideration due to laws or regulations in their jurisdiction).
    \end{itemize}


\item {\bf Broader impacts}
    \item[] Question: Does the paper discuss both potential positive societal impacts and negative societal impacts of the work performed?
    \item[] Answer: \answerYes{}
    \item[] Justification: Section~\ref{sec:disc} (Broader impact) discusses benefits for leaderboard designers and the risk that the analysis is read as a manipulation recipe.
    \item[] Guidelines:
    \begin{itemize}
        \item The answer \answerNA{} means that there is no societal impact of the work performed.
        \item If the authors answer \answerNA{} or \answerNo, they should explain why their work has no societal impact or why the paper does not address societal impact.
        \item Examples of negative societal impacts include potential malicious or unintended uses (e.g., disinformation, generating fake profiles, surveillance), fairness considerations (e.g., deployment of technologies that could make decisions that unfairly impact specific groups), privacy considerations, and security considerations.
        \item The conference expects that many papers will be foundational research and not tied to particular applications, let alone deployments. However, if there is a direct path to any negative applications, the authors should point it out. For example, it is legitimate to point out that an improvement in the quality of generative models could be used to generate Deepfakes for disinformation. On the other hand, it is not needed to point out that a generic algorithm for optimizing neural networks could enable people to train models that generate Deepfakes faster.
        \item The authors should consider possible harms that could arise when the technology is being used as intended and functioning correctly, harms that could arise when the technology is being used as intended but gives incorrect results, and harms following from (intentional or unintentional) misuse of the technology.
        \item If there are negative societal impacts, the authors could also discuss possible mitigation strategies (e.g., gated release of models, providing defenses in addition to attacks, mechanisms for monitoring misuse, mechanisms to monitor how a system learns from feedback over time, improving the efficiency and accessibility of ML).
    \end{itemize}
    
\item {\bf Safeguards}
    \item[] Question: Does the paper describe safeguards that have been put in place for responsible release of data or models that have a high risk for misuse (e.g., pre-trained language models, image generators, or scraped datasets)?
    \item[] Answer: \answerNA{}
    \item[] Justification: No models or data with high misuse risk are released; only derived score tables and code.
    \item[] Guidelines:
    \begin{itemize}
        \item The answer \answerNA{} means that the paper poses no such risks.
        \item Released models that have a high risk for misuse or dual-use should be released with necessary safeguards to allow for controlled use of the model, for example by requiring that users adhere to usage guidelines or restrictions to access the model or implementing safety filters. 
        \item Datasets that have been scraped from the Internet could pose safety risks. The authors should describe how they avoided releasing unsafe images.
        \item We recognize that providing effective safeguards is challenging, and many papers do not require this, but we encourage authors to take this into account and make a best faith effort.
    \end{itemize}

\item {\bf Licenses for existing assets}
    \item[] Question: Are the creators or original owners of assets (e.g., code, data, models), used in the paper, properly credited and are the license and terms of use explicitly mentioned and properly respected?
    \item[] Answer: \answerYes{}
    \item[] Justification: Sources and licence status are listed in \texttt{REPRODUCE.md}: neither leaderboard declares a licence for its score data, HELM code is Apache-2.0; we say so and release only derived tables.
    \item[] Guidelines:
    \begin{itemize}
        \item The answer \answerNA{} means that the paper does not use existing assets.
        \item The authors should cite the original paper that produced the code package or dataset.
        \item The authors should state which version of the asset is used and, if possible, include a URL.
        \item The name of the license (e.g., CC-BY 4.0) should be included for each asset.
        \item For scraped data from a particular source (e.g., website), the copyright and terms of service of that source should be provided.
        \item If assets are released, the license, copyright information, and terms of use in the package should be provided. For popular datasets, \url{paperswithcode.com/datasets} has curated licenses for some datasets. Their licensing guide can help determine the license of a dataset.
        \item For existing datasets that are re-packaged, both the original license and the license of the derived asset (if it has changed) should be provided.
        \item If this information is not available online, the authors are encouraged to reach out to the asset's creators.
    \end{itemize}

\item {\bf New assets}
    \item[] Question: Are new assets introduced in the paper well documented and is the documentation provided alongside the assets?
    \item[] Answer: \answerYes{}
    \item[] Justification: The released code and derived score tables are documented in \texttt{REPRODUCE.md}; no other new assets are released.
    \item[] Guidelines:
    \begin{itemize}
        \item The answer \answerNA{} means that the paper does not release new assets.
        \item Researchers should communicate the details of the dataset\slash code\slash model as part of their submissions via structured templates. This includes details about training, license, limitations, etc. 
        \item The paper should discuss whether and how consent was obtained from people whose asset is used.
        \item At submission time, remember to anonymize your assets (if applicable). You can either create an anonymized URL or include an anonymized zip file.
    \end{itemize}

\item {\bf Crowdsourcing and research with human subjects}
    \item[] Question: For crowdsourcing experiments and research with human subjects, does the paper include the full text of instructions given to participants and screenshots, if applicable, as well as details about compensation (if any)? 
    \item[] Answer: \answerNA{}
    \item[] Justification: The paper involves no crowdsourcing or human subjects.
    \item[] Guidelines:
    \begin{itemize}
        \item The answer \answerNA{} means that the paper does not involve crowdsourcing nor research with human subjects.
        \item Including this information in the supplemental material is fine, but if the main contribution of the paper involves human subjects, then as much detail as possible should be included in the main paper. 
        \item According to the NeurIPS Code of Ethics, workers involved in data collection, curation, or other labor should be paid at least the minimum wage in the country of the data collector. 
    \end{itemize}

\item {\bf Institutional review board (IRB) approvals or equivalent for research with human subjects}
    \item[] Question: Does the paper describe potential risks incurred by study participants, whether such risks were disclosed to the subjects, and whether Institutional Review Board (IRB) approvals (or an equivalent approval/review based on the requirements of your country or institution) were obtained?
    \item[] Answer: \answerNA{}
    \item[] Justification: The paper involves no human subjects, so no IRB approval applies.
    \item[] Guidelines:
    \begin{itemize}
        \item The answer \answerNA{} means that the paper does not involve crowdsourcing nor research with human subjects.
        \item Depending on the country in which research is conducted, IRB approval (or equivalent) may be required for any human subjects research. If you obtained IRB approval, you should clearly state this in the paper. 
        \item We recognize that the procedures for this may vary significantly between institutions and locations, and we expect authors to adhere to the NeurIPS Code of Ethics and the guidelines for their institution. 
        \item For initial submissions, do not include any information that would break anonymity (if applicable), such as the institution conducting the review.
    \end{itemize}

\item {\bf Declaration of LLM usage}
    \item[] Question: Does the paper describe the usage of LLMs if it is an important, original, or non-standard component of the core methods in this research? Note that if the LLM is used only for writing, editing, or formatting purposes and does \emph{not} impact the core methodology, scientific rigor, or originality of the research, declaration is not required.
    %this research? 
    \item[] Answer: \answerNA{}
    \item[] Justification: No LLM is part of the method; LLM assistance was used only for writing and coding assistance.
    \item[] Guidelines:
    \begin{itemize}
        \item The answer \answerNA{} means that the core method development in this research does not involve LLMs as any important, original, or non-standard components.
        \item Please refer to our LLM policy in the NeurIPS handbook for what should or should not be described.
    \end{itemize}

\end{enumerate}

\end{document}
